\documentclass[11pt]{article}
\usepackage[margin=1in]{geometry}
\usepackage{amsmath, amssymb, enumitem}
\setlist{noitemsep}
\usepackage{markdown}

\begin{document}

    \begin{center}
        \Large\textbf{CSCE 322 --- Worksheet 4}\\[2mm]
        \large\textbf{Lambda Calculus \& Operational Semantics}
    \end{center}

    \bigskip

    \noindent\textbf{Instructions:}
    \begin{itemize}
        \item For Part 1, explicitly show your steps for $\beta$-reduction and $\alpha$-renaming.
        \item For Part 2, use the big-step operational semantics rules from class. Write your derivations using the fraction style (premises over consequent).
        \item Assume the standard rules for numbers, booleans, variables, arithmetic, comparisons, conditionals, and \texttt{var}.
    \end{itemize}

    \newpage

    \section*{Part 1: Lambda Calculus}

    Evaluate the following lambda calculus expressions to normal form. Be careful with free and bound variables, and apply $\alpha$-renaming where necessary to avoid variable capture.

    \begin{enumerate}
        \item \textbf{Multiple Free Variables in Argument}
        \[
            (\lambda x.\ \lambda y.\ x\ (\lambda z.\ x\ y\ z))\ (\lambda y.\ \lambda x.\ y\ z\ x)
        \]

        \begin{markdown}
            use a IDE to view this because seeing the "()" without highlighting is difficult.

            rephrase into a duncan formatted lambda calculus question
            ( \lx.( \ly.( x ( \lz. ( x y z ) ) ) ) ) ( \ly. ( \lx. ( y z x ) ) )

            do \a renaming
            ( \lx2.( \ly.( x2 ( \lz1. ( x2 y z1 ) ) ) ) ) ( \ly1. ( \lx1. ( y1 z x1 ) ) )

            plug in right side for x2
            _____        __                             _______________________________
            ( \lx2.( \ly.( x2 ( \lz1. ( x2 y z1 ) ) ) ) ) ( \ly1. ( \lx1. ( y1 z x1 ) ) )
            -->
            ( \ly.( ( \ly1. ( \lx1. ( y1 z x1 ) ) ) ( \lz1.( ( \ly1.( \lx1.( y1 z x1 ) ) )
            y z1 ) ) ) )

            cannot plug var into \ly. no var present. go to \ly1.
            plug in \lz1 for y1
            ( \ly. ( \lx1. ( ( \lz1. ( ( \ly1.( \lx1.( y1 z x1 ) ) ) y z1 ) ) z x1 ) ) )

            cannot plug var into \lx1. no var present go to \lz1.
            plug in z for z1
            _____                                   __     _
            ( \ly. ( \lx1. ( ( \lz1. ( ( \ly1.( \lx1.( y1 z x1 ) ) ) y z1 ) ) z x1 ) ) )
            -->
            ( \ly. ( \lx1. ( ( ( \ly1. ( \lx1. ( y1 z x1 ) ) ) y z ) x1 ) ) )

            plug in y for y1
            _____           __            _
            ( \ly. ( \lx1. ( ( ( \ly1. ( \lx1. ( y1 z x1 ) ) ) y z ) x1 ) ) )
            -->
            ( \ly. ( \lx1. ( ( ( \lx1. ( y z x1 ) ) z ) x1 ) ) )

            plug in z for \lx1.
            _____       __     _
            ( \ly. ( \lx1. ( ( ( \lx1. ( y z x1 ) ) z ) x1 ) ) )
            -->
            ( \ly. ( \lx1. ( ( y z z ) x1 ) ) )

            this problem has more memes that make it harder than traditional
            plug and play lambda calculus problems
            for here, you have to see if substitution is possible.
            final answer:
            ( \ly. ( \lx1. ( ( y z z ) x1 ) ) )
        \end{markdown}

        \newpage

        \item \textbf{Shadowing and Scope}
        \[
            ((\lambda x.\ \lambda y.\ x\ (\lambda y.\ x\ y))\ (\lambda z.\ y\ z))\ x
        \]

        \begin{markdown}
            convert this into a duncan lambda calculus question
            ( ( \lx. ( \ly. ( x ( \ly. ( x y ) ) ) ) ) ( \lz. ( y z ) ) ) x

            \a renaming time :-DDD
            ( ( \lx1. ( \ly. ( x1 ( \ly1. ( x1 y1 ) ) ) ) ) ( \lz. ( y z ) ) ) x

            plug in \lz. and it's buddies for x1
            _____          __           __              __________________
            ( ( \lx1. ( \ly. ( x1 ( \ly1. ( x1 y1 ) ) ) ) ) ( \lz. ( y z ) ) ) x
            -->
            ( \ly. ( ( \lz. ( y z ) ) ( \ly1. ( ( \lz. ( y z ) ) y1 ) ) ) ) x

            plug in x for y
            ____            _                          _                  _
            ( \ly. ( ( \lz. ( y z ) ) ( \ly1. ( ( \lz. ( y z ) ) y1 ) ) ) ) x
            -->
            ( \lz. ( x z ) ) ( \ly1. ( ( \lz. ( x z ) ) y1 ) )

            plug in \ly1 and it's buddies for z
            ____     _     _________________________________
            ( \lz. ( x z ) ) ( \ly1. ( ( \lz. ( x z ) ) y1 ) )
            -->
            ( x ( \ly1. ( ( \lz. ( x z ) ) y1 ) ) )

            cannot plug var into x, not a \l function.
            cannot plug var into y1, no var available.
            plug in y1 for z
            ____     _     __
            ( x ( \ly1. ( ( \lz. ( x z ) ) y1 ) ) )
            -->
            ( x ( \ly1. ( x y1 ) ) )

            cannot plug var in for y1, no var available.
            seems like that's the end.
            ( x ( \ly1. ( x y1 ) ) )
        \end{markdown}

        \newpage

        \vspace{1in}

        \item \textbf{Function Composition with Free Variables}
        \[
            (\lambda f.\ \lambda g.\ \lambda x.\ f\ (g\ x))\ (\lambda x.\ y\ x)\ (\lambda y.\ x\ y)
        \]

        \begin{markdown}
            convert to duncan lambda calculus question
            ( \lf. ( \lg. ( \lx. ( f ( g x ) ) ) ) ) ( \lx. ( y x ) ) ( \ly. ( x y ) )

            \a renaming :DDDD
            ( \lf. ( \lg. ( \lx2. ( f ( g x2 ) ) ) ) ) ( \lx1. ( y x1 ) ) ( \ly1. ( x y1 ) )

            plug in x1 and it's buddies for f
            ____                  _                  __________________
            ( \lf. ( \lg. ( \lx2. ( f ( g x2 ) ) ) ) ) ( \lx1. ( y x1 ) ) ( \ly1. ( x y1 ) )
            -->
            ( \lg. ( \lx2. ( ( \lx1. ( y x1 ) ) ( g x2 ) ) ) ) ( \ly1. ( x y1 ) )

            plug in y1 for g
            ____                                _            __________________
            ( \lg. ( \lx2. ( ( \lx1. ( y x1 ) ) ( g x2 ) ) ) ) ( \ly1. ( x y1 ) )
            -->
            ( \lx2. ( ( \lx1. ( y x1 ) ) ( ( \ly1. ( x y1 ) ) x2 ) ) )

            plug in y1 and it's buddies for x1
            _____     __     _________________________
            ( \lx2. ( ( \lx1. ( y x1 ) ) ( ( \ly1. ( x y1 ) ) x2 ) ) )
            -->
            ( \lx2. ( y ( ( \ly1. ( x y1 ) ) x2 ) ) )

            cannot plug var in for y, no \ly present
            plug in x2 for y1
            _____     __     __
            ( \lx2. ( y ( ( \ly1. ( x y1 ) ) x2 ) ) )
            -->
            ( \lx2. ( y ( x x2 ) ) )

            cannot reduce further.
            final answer:
            ( \lx2. ( y ( x x2 ) ) )
        \end{markdown}

        \vspace{1in}
    \end{enumerate}

    \newpage

    \section*{Part 2: Operational Semantics}

    \begin{enumerate}
        \item \textbf{Deep Shadowing \& Nested Evaluation} \\
        Give a full big-step derivation for the following expression. Assume an arbitrary initial environment $\rho$.
        \[
            \langle (var\ (x\ 5)\ (var\ (y\ (if\ (gt\ x\ 3)\ (+\ x\ 2)\ 0))\ (var\ (x\ (+\ y\ 1))\ (+\ x\ y)))),\ \rho \rangle
        \]

        \begin{markdown}
            I'll start with the evaluation 1st, then the "formal definition"
            crap that looks like a bacterium crawling on your screen.

            convert this into a duncan friendly question. That's my name btw.
            ( ( var ( x 5 ) ( var ( y ( if ( gt x 3 ) ( + x 2 ) 0 ) )
            ( var ( x ( + y 1 ) ) ( + x y ) ) ) ) , p )


        \end{markdown}

        \vspace{3in}

        \item \textbf{Late Stuck State} \\
        Give a big-step derivation showing exactly where the following expression gets stuck. Explain why the derivation cannot continue.
        \[
            \langle (if\ (eq\ 1\ 1)\ (var\ (x\ true)\ (if\ x\ (+\ x\ 1)\ 5))\ 0),\ \rho \rangle
        \]

        \vspace{3in}

        \item \textbf{Complex Conditionals \& Initial Environment} \\
        Let $\rho_0 = \{ a \mapsto 10, b \mapsto 5 \}$. Evaluate the following expression:
        \[
            \langle (if\ (and\ (lt\ b\ a)\ (eq\ (-\ a\ b)\ 5))\ (var\ (a\ b)\ (+\ a\ b))\ 0),\ \rho_0 \rangle
        \]

    \end{enumerate}

\end{document}
